\Article{368}
We say that a congruence $S\equiv0$ has \emph{equal} roots, or more generally equal divisors, if it is divisible by a power of some function.

\SourceLeaf{246}
Whether a given congruence has equal divisors is decided in the same way as in the theory of equations. Put
\[
X\equiv\xi^mP
\]
Clearly,
\[
\frac{\mathrm{d}X}{\mathrm{d}x}\equiv\xi^{m-1}\left(mP\frac{\mathrm{d}\xi}{\mathrm{d}x}+\xi\frac{\mathrm{d}P}{\mathrm{d}x}\right)
\]
so $\frac{\mathrm{d}X}{\mathrm{d}x}$ is divisible by $\xi^{m-1}$. More generally, let
\[
X\equiv A^aB^bC^c\text{ etc.}
\]
where $A,B,C$, etc., denote distinct prime functions. Then
\[
\frac{\mathrm{d}X}{\mathrm{d}x}\equiv X\left(\frac{a\,\mathrm{d}A}{A\,\mathrm{d}x}+\frac{b\,\mathrm{d}B}{B\,\mathrm{d}x}+\frac{c\,\mathrm{d}C}{C\,\mathrm{d}x}+\text{etc.}\right)
\]
It follows that, unless one of the numbers $a,b,c$, etc., is divisible by the modulus, $\frac{\mathrm{d}X}{\mathrm{d}x}$ is divisible by $A^{a-1}B^{b-1}C^{c-1}$, etc., but not by $A^a,B^b,C^c$, etc. Hence the following

\textsc{Theorem.} \emph{If the common divisor of greatest degree of the functions $X$ and $\frac{\mathrm{d}X}{\mathrm{d}x}$ is $\xi$, all the prime factors of $\xi$ will also occur in $X$, each with a multiplicity $+1$ greater than in $\xi$. Thus, if $X$ and $\frac{\mathrm{d}X}{\mathrm{d}x}$ are coprime functions, $X$ will have no equal factors.}

\Needspace{17\baselineskip}\Article{369}
Example I. We ask whether the function
\[
x^5+3x^4-6x^3+3x-4\ldots(X)
\]
has equal divisors modulo $17$. We have
\[
\frac{\mathrm{d}X}{\mathrm{d}x}\equiv5x^4-5x^3-xx+3
\]
From this we find that the functions $X$ and $\frac{\mathrm{d}X}{\mathrm{d}x}$ are coprime, so $X$ has no equal divisors.

Example II. Let
\[
X\equiv x^5+6x^4-3x^3-4xx+2x-3(\mathrm{mod.}\,13)
\]
Then
\[
\frac{\mathrm{d}X}{\mathrm{d}x}\equiv5x^4-2x^3+4xx+5x+2
\]
The greatest common divisor of the functions $X,\frac{\mathrm{d}X}{\mathrm{d}x}$ is $\equiv5xx+7x+7$, or,
\SourceLeaf{247}
on multiplication by $8$: $xx+4x+4$. Since this divisor is $\equiv(x+2)^2$, the function $X$ is divisible by $(x+2)^3$, and the quotient (which is $xx+11$) contains no further double divisor.

\Article{370.\ 371}
If, by the preceding sections, the function $X$ has been expressed as $A^aB^bC^c$, etc., with $A,B,C$, etc., coprime and the numbers $a,b,c$, etc., unequal, the resolution can be carried further. Let $X$, then, be a function containing no further equal divisors. We saw above that $x^p-x$ is the product of all prime functions of degree one. Let $\xi$ be the common divisor of greatest degree of the functions $X$ and $x^p-x$; then $\xi$ is the product of all divisors of $X$ of degree one, and $\frac X\xi$ has no further divisors of this kind. If, however, the functions $X$ and $x^p-x$ are found to be coprime, $X$ will have no divisor of degree one, and hence the congruence $X\equiv0$ will have no real roots. Further, since $x^{pp}-x$ is the product of all prime functions of degrees two and one, the common divisor of greatest degree of the functions $x^{pp}-x$ and $\frac X\xi$, namely $\xi'$, contains all divisors of $X$ of degree two. Continuing in this way, we see that $X$ is thus resolved into factors $\xi,\xi',\xi''$, etc., which contain all divisors of degrees one, two, three, etc., respectively.

\Article{372}
But if a product of several prime functions of the same degree is given, the individual functions must be found by trial. This problem is closely analogous to that of finding the factors of composite numbers. Here, however, it is determined a priori whether the given function can still be split into factors. Since here too the number of all possible factors is finite, we can use an aid similar to the one above. But we do not wish to dwell on this matter, for an experienced calculator who has thoroughly grasped the principles will readily find particular devices when needed.

We proceed to another topic, namely the consideration of congruences when the modulus is not a prime number, as we have always assumed up to now. Here the case where the modulus is a power of a prime number especially deserves attention, both in itself and because it is needed to remove certain doubts (\textsection.\textsection.\ .\ .).
